% ===================================================================
%  Section X: Related Work
% ===================================================================

\paragraph{Offline RL and DQN.}
DQN combines value-function approximation, experience replay, and a target
network for temporal-difference learning~\cite{mnih2015dqn,
watkins1992qlearning}. Double DQN~\cite{hasselt2016double} separates
next-action selection from target-network value evaluation to reduce
overestimation. Offline RL studies the fixed-dataset
setting~\cite{levine2020offline}, where value learning is brittle under
distribution shift~\cite{fujimoto2019offpolicy}; Conservative Q-Learning
(CQL)~\cite{kumar2020cql} addresses this with value regularization. Our work
does not improve the RL algorithm: it makes selected DQN-style replay and TD
computations \emph{verifiable} over committed data.

\paragraph{Proof-of-learning and proof-of-training.}
Proof-of-Learning~\cite{jia2021proofoflearning} introduced the problem of
proving that a model resulted from a claimed training procedure. The approach
records intermediate checkpoints and verifies training consistency, but it is
not zero-knowledge and does not provide cryptographic soundness guarantees.
zkDL~\cite{zkdl2023} targets zero-knowledge proofs for deep learning training
including neural nonlinearities and backpropagation, using custom arithmetic
circuits. Kaizen~\cite{kaizen2024} gives a zero-knowledge proof-of-training
construction for DNNs based on per-iteration proofs and recursive composition,
achieving end-to-end training verification for networks up to VGG-11 scale.
zkPoT~\cite{zkpotvicinity2025} explores an alternative approach through optimum vicinity, proving that a
trained model lies near the optimum rather than re-executing every gradient
step. Tan et al. also name the prover's freedom over training randomness as a
bias vector, and answer it by leaving the trajectory unproved: optimum vicinity
bounds the distance to the optimum of a convex problem. That answer does not
transfer here, because offline $Q$-learning is not convex and offers no optimum
to be near, so we bind the randomness rather than bypass it --- the sampler seed
is derived from the dataset commitment inside the relation
(\Cref{thm:sampler-binding}). These systems motivate our goal but do not
directly encode offline-RL-specific semantics such as terminal Bellman branches, replay membership, Double-DQN
argmax, target-network synchronization, and TD losses.

Work published since has filled this space by model class rather than by
learning paradigm, and that is the pattern this work reads as an opening.
ZKBoost~\cite{zkboost2026} gives the first zero-knowledge proof of training for
gradient-boosted trees, using a fixed-point XGBoost that stays within $1\%$ of
the floating-point model --- the same discipline our relations follow, arrived
at independently. SUMMER~\cite{summer2026} gives the first end-to-end zkPoT for
recurrent training, folding per-step commitments with Merkle trees so that
proof size and verification time are independent of the iteration count --- the
same structural lever our aggregation pulls. VeriLoRA~\cite{verilora2026}
carries verifiable fine-tuning to $13$B-parameter language models by
restricting the proved computation to the low-rank adapters. Trees, recurrences
and adapters have each been covered; sequential decision-making has not, and
the reason is not model size --- our networks carry fewer than a thousand
parameters --- but the three obligations of \Cref{sec:introduction}.

\paragraph{Verifiable and zero-knowledge machine learning.}
Verifiable ML has been studied extensively for inference integrity.
SafetyNets~\cite{safetynets2017} builds proofs for neural-network inference on
untrusted clouds. zkCNN~\cite{zkcnn2021} proves CNN predictions and accuracy
while hiding private information. VeriML~\cite{veriml2021} provides integrity
and payment assurances for outsourced ML tasks. vCNN~\cite{vcnn2024} implements
verifiable convolutional neural networks using zk-SNARKs and was published in
IEEE TDSC. Recent surveys organize this space into verifiable inference,
testing, and training~\cite{zkmlsurvey2026}. Our relation is different: the
first-class computation is a Bellman TD claim over committed replay data rather
than a supervised prediction.

\paragraph{ZK for federated learning.}
Zero-knowledge proofs have been applied to verify gradient aggregation in
federated learning~\cite{zkgradient2025}, where a ZK proof ensures that the
aggregator correctly combined client updates. While this addresses a different
threat model (malicious aggregator vs.\ malicious trainer), the gradient
verification component shares technical similarities with our backpropagation
and SGD update proofs.

\paragraph{zkVMs and proof backends.}
General-purpose zkVMs trade circuit specialization for implementation speed.
SP1~\cite{succinct2024sp1} is a Rust-oriented RISC-V zkVM using STARK-based
proofs; we chose it because our relations combine hashing, Merkle paths, schema
checks, integer arithmetic, and branching. RISC Zero~\cite{risczero2023} is
another zkVM in the same design space. Both support recursive proof
composition, and we use it in both available forms: a proof-manifest chain that
binds externally verified child manifests, and aggregation in which the guest
verifies each child proof cryptographically in-circuit. On-chain verification
can be achieved through
Groth16~\cite{groth2016zksnark} or PLONK~\cite{gabizon2019plonk} SNARK
compression.

\paragraph{Verifiable RL, and a name collision worth separating.}
The phrase now carries two unrelated meanings, and only one of them is ours.
In current usage ``verifiable reinforcement learning'' most often means RL with
verifiable rewards, where a checker --- unit tests, an exact-answer comparison,
a proof assistant --- supplies the training signal. Nothing is proved to a
third party there; the verifier sits inside the training loop and the model
owner is the one who benefits from its verdict. Our verifier sits outside, is
adversarial to the prover, and learns nothing beyond the statement. The older
line of verifiable RL is likewise about \emph{policies} rather than training:
programmatic-policy approaches~\cite{verma2019verifiable} search for policies
checkable by symbolic methods, which certifies what a policy will do, not how
it came to be. Our work targets a different layer:
cryptographic verification of the replay membership, Bellman arithmetic,
model-grounded TD evidence, and update evidence inside offline value learning.
To the best of our knowledge, this is the first system to provide
zero-knowledge proof-backed verification for selected offline-RL training
relations over committed replay data.

\paragraph{Data provenance for AI.}
The importance of training-data provenance is increasingly
recognized~\cite{dataauth2024, aiprovenance2025}. Frameworks for cryptographic
verifiability of end-to-end AI pipelines propose Merkle commitments and content
identifiers for dataset integrity. ZKPROV~\cite{zkprov2025} binds a language
model's responses to datasets certified by their owners, proving provenance
without revealing dataset contents. We read it as complementary to
proof-of-training rather than a substitute: it establishes provenance of the
data, not correctness of the computation over it. Our commitment sits on the same side of
that line but is consumed on the other: the Merkle root that
\Cref{alg:dataset-commitment} produces is an input the training relation reads,
so the sampler, the membership check and the aggregate root all bind to it. A
provenance claim and a training claim that name the same root are the two
halves a regulator needs, and joining them is cheaper than proving either.

\paragraph{Verifiable computation.}
The broader verifiable computation literature~\cite{walfish2015verifiable}
provides the theoretical foundations for our approach. Our system applies these
foundations to a specific domain (offline RL training) and demonstrates that
general-purpose zkVM technology can be used to verify domain-specific machine
learning computations at a practical proof cost.
